\begin{lemma}[Three-letter lifting with negative outer signs]\label{lem:three-letter-negative-lift}
With \(L=S_g\) and \(B\) the short-band indicator of Section~\ref{sec:inert-row-table},
for the chronological signs \(-,+,-\), one can choose a three-letter correction in the short-arc frames with exponent \(e_3\) satisfying
\[
 e_3\le22L+1+\eta B
\]
in the given DVR normalization.
\end{lemma}

\begin{proof}
\label{par:short-three-levels}
Put
\[
 u=s_{c-2},\ v=s_{c-1},\ z=s_c,
\]
so \(u\le w\)
and \(v,z\in[b_0,b_0+\eta]\). Here \(40L\le b_0/4\).
At a partial arc use level numbers relative to its own endpoint.

\proofstep{Start with the middle positive letter and retain the two negative neighbors.}
First consider signs \(-,+,-\) in chronological order.
If \(p(u+S+\eta)\le v\), work at the middle letter and
write \(A=A_{c-1}\), \(y=u+S+\eta\).
As in the two-letter coefficient extraction use the coordinate
\(P_1^{u+S}/q^v\) in direction \(d\). At this intermediate
point the scalar image bound at a cut is unavailable. Instead
lift \(X^A\) times the realized vector. After the two horizontal
comparisons and multiplication by \(q^{v+S}P_1^\eta\), the
scalar on the target, in the variable \(q\), is
\[
 F=p^{2\sigma}(q-p)^A q^S P_1^y .
\]
Below degree \(v+S-\eta\) every source coefficient lies in
\(\mathcal F^v\) by the gap. To recover an integral direction
it suffices that some such coefficient of \(F\) divides a small
power \(p^t\) in \(O\).

Take
\[
 t=1+\left\lfloor\tfrac{11}{10}
 (2\sigma+(\eta+A)/p+2+\gamma_h)\right\rfloor .
\]
Then \(2\sigma<t\le10L\), \(A+t+\eta\le v\).
For the first bound use \(\sigma\le3L\) and
\(A\le (h+\eta)/(p-1)+1\), giving
\[
 (\eta+A)/p\le 1.1(h+\eta)/p^2+(\eta+1)/p.
\]

Substitution in the floor bound gives \(t\le10L\).
For the second, \(A\le (b_0+2\eta+S)/(p-1)+1\),
\(10L\le b_0/16\) and \(20\eta,20S\le b_0\).
If all the indicated coefficients were divisible by \(p^t\),
put \(k=t-2\sigma\). Modulo \(p^k\) one has
\(q^{A+k-1}\in((q-p)^A)\) by the binomial expansion.
Cancel \(p^{2\sigma}\), remove the initial power of \(q\),
and use this divisibility. The result would be
\[
 P_1^y\in(p^k,q^{v-\eta-(A+t-1)}).
\]
The second content calculation, before rounding, gives
\[
 k\le\max\{0,\ y-(v-\eta-A-t+1)/p+2+\log_p y\}.
\]
Indeed using the index \(\min(y,\lfloor(C-1)/p\rfloor)\)
costs at most \(y-C/p+1+1/p+\log_p y\) when positive.

Since \(py\le v\) and \(p/(p-1)\le11/10\), the definition
of \(t\) contradicts the last display. Choose a coefficient
of valuation less than \(t\), and multiply its source vector by
the integral quotient of \(p^t\) by that coefficient.
Moving one step forward gives the positive-level preimage.
The sole pole of \(f\) is at this level, so the second lifting
rule gives cost
\[
 t+\sigma\le13L.
\]


Suppose next that \(v<p(u+S+\eta)\). The down-level constants
in \eqref{eq:short-scalar} can be different. Eliminate the one
at level two first. The denominator there is \(P_1^v\), the
modulus is \(P_2^{u+S}\). Put \(y=u+S\). Since
\(P_2=p\varepsilon_1+U_1P_1^p\) with unit coefficients, the
first content calculation (using \(P_1,P_2\) in place of
\(q,P_1\)) gives a splitting
\[
 t_2f=x_2+f_2,\quad D f_2\in P_2^{y}\mathcal R,\quad
 D(p^{k_0}f_2-x_0)\in q^{z+S}\mathcal R,\qquad D=P_1^v .
\]

If \(v<py\), the content of \(P_1^v\) modulo \(P_2^y\)
is zero by reduction and comparison of orders. Otherwise use
the ceiling multiple \(p\lceil v/p\rceil\) in its place.
For \(py\le v<p(y+\eta)\), the expansion in \(p,P_1^p\) of \(P_2^y\) costs
at most \(\eta+2+\gamma_h\), and the small letter belongs to
the band. Thus \(v_p(t_2)\le(\eta+2+\gamma_h)B\).

Write \(C=z+S\). For \(x_0\ne0\), the second congruence gives
\begin{equation}\label{eq:short-last-constant}
 v_p(x_0)-k_0\ge y-v+C/p-C/p^2-2-\gamma_h .
\end{equation}
Indeed \(P_2^y\) has content at least \(y-C/p^2\) modulo
\(q^C\). Expand first in \(p,P_1^p\), then in \(p,q^p\):
a term with less scalar valuation has \(q\)-degree at least
\(C\). The content of \(P_1^v\) modulo \(q^C\) is at most
\(v-C/p+2+\gamma_h\) by the second calculation, since
\(C\le pv\). Comparing contents of
\[
 x_0P_1^v\equiv p^{k_0}P_2^y U\bmod q^C
\]
proves the claim.

Here is how to use this scalar inequality with the middle
filtered lattice. Transport the truncated horizontal column to
the middle point, paying \(\sigma\) at the first small letter,
and denote its integral representative by \(\lambda\). It
represents \(p^\sigma d\) there.

Set \(s=P_1^y/q^v\) at this
point. The vector \(s\lambda\) lies in the target Wach lattice.
The level test uses \(d\in\mathcal F^v\) at the middle and
the zero of order \(y\) at the initial negative letter.
In addition
\[
 \gamma(s)(\gamma-1)\lambda\in Xq^{a_*}N_{c-1}(T_2)
\]
where \(a_*=A_{c-1}+r_{c-1}\) and \(r_i\) is the least
scalar Wach image exponent at \(i\). Indeed the one-step
horizontal correction has defect in \(X^M q^{2G}\) in the
first Frobenius span. The next transport gives factors
\(X^M q^{(p-1)M}P_1^{2G}\); after multiplication by
\(\gamma(s)\) these pay \(Xq^{a_*}\) and the signed height losses
\(q^\eta P_1^{u+S}\), since \((p-1)M\ge v+a_*+\eta\).

The powers required in the ordinary parameter and the scalar
give the \(q^{a_*}\) image by the binomial theorem.
The two-filter argument of Paragraph~\ref{par:thick-filtered-comparison}
therefore lifts \(p^E s\lambda\) integrally, where
\[
 E=\left\lfloor\frac{3a_*+S+2\eta-2}{p-1}
               +\log_p(v+\eta-1)\right\rfloor_{\!+}\le8L .
\]
For this application the lower filter has indices \([v-a_*+1,v]\)
and the other has indices in \([v+1,v+S+a_*]\cup[1-\eta,\eta+a_*]\).
The factors adjacent to \(j\) in the first two intervals
give factorials of combined length \(2a_*+S-1\).

In the remaining
interval of length \(2\eta+a_*\), isolate a term of largest
valuation; the others are bounded by factorials of combined
length \(2\eta+a_*-1\). The isolated valuation is at most
\(\log_p(v+\eta-1)\). If \(a_*=0\) the map is already onto
at this letter. The global ordinary width bound \(H\le h+\eta\)
and degree lowering give
\[
 a_*\le H(2p-1)/(p-1)^2+\lceil1+\log_p(1+H)\rceil+1.
\]

These imply \(2a_*+\eta+1\le v\) under \(40L\le b_0/4\)
and give the displayed estimate for \(E\), by \(p\ge11\) and
\(H\le b_0+2\eta+S\). In particular
\[
 v+a_*+\eta\le2v\le2h<(p-1)M,
\]
as used in the defect estimate. The separation gives the
disjointness required by interpolation. \label{cor:inert-short-three-1}

Let \(n_1\) be this lift. Take its constant in a horizontal
basis at \(c-1\) and transport it to \(c\) with precision
\(Xq^C\). This gives \(y_1\) in the last Frobenius span,
with \((\gamma-1)y_1\in XN_c(T_1)\) and \(q^C y_1\in N_c(T_1)\).
Use the truncation order
\[
 R_0=\lceil C/(p-1)\rceil
\]
and
take \(\sigma_1=R_0\). In positive degree \(r<R_0\) the
denominator costs \(r+v_p(r!)\). The transported monomial is
a unit times \(X^r q^{(p-1)r}\). Keep one \(X\) and expand
the others as \((q-p)^{r-1}\).

Modulo \(Xq^C\) the surviving
terms contain at least \(\max(0,pr-C)\) powers of \(p\):
a term of extra degree \(j\) must satisfy \(j+(p-1)r<C\).
The loss is therefore at most
\[
 pr/(p-1)-\max(0,pr-C)\le C/(p-1);
\]
the maximum occurs at the crossing \(pr=C\). The transported
tail is zero to this precision. The last-level jet of \(y_1\)
is horizontal in the last Frobenius span: its difference is
zero modulo \(Xq^C\) there, by the truncation and the integral
correction. Over the coefficient field use the local logarithmic
parameter at the primitive level. On a horizontal column the
action on degree \(k\) has eigenvalue \(\chi(\gamma)^k\), so
an invariant jet has only its constant coefficient. This gives
the jet of \(F(y_1)\) from the transported constant of \(n_1\)
with multiplier \(p^{\sigma_1}\).

To compare with \(f_2\Phi\lambda\), use the common horizontal
column at the last level. At this level the transported multiplier
\(\phi(s)=P_2^y/P_1^v\) has constant term
\[
 (P_2^y/P_1^v)(0)=p^{y-v},
\]
since \(P_j(0)=p\) for \(j\ge1\) in the \(q\)-coordinate.
Thus the two scalar constants to be matched, after multiplying
by \(p^{k_0+v}\), are \(x_0p^v\) and
\(p^{k_0+\sigma_1+E+y}\). In the second expression the
common horizontal normalization cancels; its \(p^\sigma\) is
present in both comparisons.

Choose \(t_3,B'\in O\) with
\[
 t_3 x_0p^v=B' p^{k_0+\sigma_1+E+y},\qquad v_p(t_3)\le14L .
\]
For \(x_0=0\) use \((1,0)\). Otherwise take \(t_3=1\) if
\(v_p(x_0)+v\ge k_0+\sigma_1+E+y+\sigma\). In the other
case choose equality with \(B'=p^\sigma\). Equation
\eqref{eq:short-last-constant} bounds \(v_p(t_3)\) by
\[
 \sigma+E+\sigma_1-C/p+C/p^2+2+\gamma_h\le14L .
\]
In fact \(\sigma_1\le C/(p-1)+1\), \(C\le h\), and
\(1/(p(p-1))+1/p^2\le 2.1/p^2\).

For the marked map \(F\), the required Frobenius-span congruence is
\[
 t_3 f_2\Phi\lambda-B' F(y_1)\in q^C\langle\Phi N_{c-1}(T_2)\rangle.
\]

For the first term the numerator at level two supplies the
height loss from the first negative letter; equivalently
\(f_2\Phi\lambda\) is an integral multiple of
\(\Phi(s\lambda)\). For the divisibility assertion compare
the two horizontal jets after inverting \(p\) and the higher
level factors. Modulo \(q^C\) the first has constant multiplier
\(x_0/p^{k_0}\) on \(\Phi\lambda\), by the last congruence;
the horizontal jet of the second has constant multiplier
\(p^{\sigma_1+E+y-v}\). Higher level factors are units at
\(q=0\) over the fraction field, and the bounded integral
Frobenius span is free.

For the free bounded integral Frobenius span \(P\), monic division yields
\[
 P\cap q^C P[1/p]=q^C P.
\]
Thus the congruence of rational horizontal jets is integral. The negative terminal
height puts this difference in \(N_c(T_2)\). Lift \(p^R\)
times the difference there and add \(p^R B' y_1\).
The result is admissible and maps to \(p^Rt_3 f_2\Phi\lambda\);
its denominators in \(N_c(T_1)\) are cleared by \(q^C\).
Since \(\Phi\lambda\) represents \(p^\sigma d\), this
lifts \(t_3f_2\) at cost \(R+\sigma\). The summand \(t_3x_2\)
admits the constant lift at the same cost. The scalar clearing cost is
\[
 \lceil v_p(t_3t_2)\rceil=\lceil v_p(t_3)+v_p(t_2)\rceil.
\]
Adding the lifting costs gives the bound \(22L+1+\eta B\).
\end{proof}

